% sec3_7.tex — Section 3.7: Hypothesis Testing by the Likelihood-Ratio Principle

We have derived in Section~3.5 the chi-squared test statistics for the null hypothesis
$\mathrm{H}_0\colon \mathbf{a}(\boldsymbol{\delta}) = \mathbf{0}$ by the Wald principle.  This section does
the same by the likelihood-ratio (LR) principle, which is to examine the difference in the
objective function with and without the imposition of the null.  Derivation of test
statistics by the Lagrange Multiplier principle for GMM and extension to nonlinear
equations are given in Section~7.4.

In the efficient GMM estimation, the objective function is
$J(\tilde{\boldsymbol{\delta}}, \hat{\mathbf{S}}^{-1})$ for a given consistent estimate $\hat{\mathbf{S}}$ of
$\mathbf{S}$.  The \textbf{restricted efficient GMM estimator} is defined as
\begin{equation}
  \text{restricted efficient GMM:} \quad
  \tilde{\boldsymbol{\delta}}(\hat{\mathbf{S}}^{-1})
  \equiv \argmin_{\tilde{\boldsymbol{\delta}}}\;
  J(\tilde{\boldsymbol{\delta}},\, \hat{\mathbf{S}}^{-1})
  \quad \text{subject to } \mathrm{H}_0.
  \label{eq:3.7.1}
\end{equation}
The LR principle suggests that
\begin{equation}
  LR \equiv J\bigl(\tilde{\boldsymbol{\delta}}(\hat{\mathbf{S}}^{-1}),\,
  \hat{\mathbf{S}}^{-1}\bigr)
  - J\bigl(\hat{\boldsymbol{\delta}}(\hat{\mathbf{S}}^{-1}),\,
  \hat{\mathbf{S}}^{-1}\bigr)
  \label{eq:3.7.2}
\end{equation}
should be asymptotically chi-squared.  Indeed it is.

\begin{proposition}[test statistic by the LR principle]
\label{prop:3.8}
\emph{Suppose Assumptions~3.1--3.5 hold and suppose there is available a consistent estimator,
$\hat{\mathbf{S}}$, of $\mathbf{S}$ $(= \E(\mathbf{g}_i \mathbf{g}_i'))$.
Consider the null hypothesis of $\#\mathbf{a}$ restrictions $\mathrm{H}_0\colon
\mathbf{a}(\boldsymbol{\delta}) = \mathbf{0}$ such that $\mathbf{A}(\boldsymbol{\delta})$, the
$\#\mathbf{a} \times L$ matrix of first derivatives, is continuous and of full row rank.  Define two
statistics, $W$ and $LR$, by \eqref{eq:3.5.16} and \eqref{eq:3.7.2}, respectively.  Then,
under the null, the following holds:}
\begin{enumerate}[label=(\alph*)]
\item \emph{The two statistics are asymptotically equivalent in that their asymptotic
  distributions are the same (namely, $\chi^2(\#\mathbf{a})$).}

\item \emph{The two statistics are asymptotically equivalent in the stronger sense that their
  numerical difference converges in probability to zero: $LR - W \pto 0$.
  (By Lemma~2.4(a), this result is stronger than~(a).)}

\item \emph{Furthermore, if the hypothesis is linear so that the restrictions can be written as
  $\mathbf{R}\boldsymbol{\delta} = \mathbf{r}$, then the two statistics are numerically equal.}
\end{enumerate}
\end{proposition}

Proving this for the linear case, which is an algebraic result, is left as an analytical
exercise.  For a full proof, see Section~7.4.

Several comments about Proposition~3.8:

\begin{itemize}
\item The advantage of $LR$ over $W$ is invariance: the numerical value of $LR$ does not
  depend on how the nonlinear restrictions are represented by $\mathbf{a}(\cdot)$.  On the other
  hand, you have to write a nonlinear optimization computer program to find the
  restricted efficient GMM when the hypothesis is nonlinear.

\item Proposition~3.8 requires that the distance matrix $\hat{\mathbf{W}}$ satisfy the efficiency
  condition $\plim\, \hat{\mathbf{W}} = \mathbf{S}^{-1}$.  Otherwise $LR$ is not asymptotically chi-squared.
  In contrast, the Wald statistic is asymptotically chi-squared without $\hat{\mathbf{W}}$
  satisfying the efficiency condition.

\item The same estimate of $\mathbf{S}$ should be used throughout in the calculation of $LR$.
  Let $\tilde{\mathbf{S}}$ and $\bar{\mathbf{S}}$ be two different consistent estimators of $\mathbf{S}$,
  and consider the statistic
  \[
    J(\tilde{\boldsymbol{\delta}}(\bar{\mathbf{S}}^{-1}),\, \bar{\mathbf{S}}^{-1})
    - J(\hat{\boldsymbol{\delta}}(\tilde{\mathbf{S}}^{-1}),\, \tilde{\mathbf{S}}^{-1}).
  \]
  This statistic is what you end up with if you perform two separate two-step efficient
  GMMs with and without the constraint of the null; $\bar{\mathbf{S}}$ is the estimate of $\mathbf{S}$
  from the first step with the constraint, while $\tilde{\mathbf{S}}$ is from the first step without the
  constraint.  The statistic is asymptotically equivalent to $LR$ (in that the difference
  between the two converges in probability to zero), but in finite samples it may
  be negative.  Having the same estimate of $\mathbf{S}$ throughout ensures the nonnegativity
  of the statistic in finite samples.  Researchers usually use the estimate from
  unconstrained estimation ($\tilde{\mathbf{S}}$) here, but the estimate from constrained estimation
  is also valid because it is consistent for $\mathbf{S}$ under the null.

\item Part~(b) of the proposition (the asymptotic equivalence in a stronger sense)
  means that if the sample size is large enough and the hypothesis is true, then the
  outcome (not just the probability of rejection or acceptance) of the test based
  on $LR$ will be the same as that based on $W$ because the probability that the two
  statistics differ numerically by an even tiny amount is zero in sufficiently large
  samples.

\item For the numerical equivalence $LR = W$ for the linear case to hold, the same
  $\hat{\mathbf{S}}$ must be used throughout to calculate not only $LR$ but also the Wald statistic.
  Otherwise $LR$ and $W$ are only asymptotically equivalent.
\end{itemize}

%% ─── The LR Statistic for the Regression Model ────────────────────────────
\subsection*{The $LR$ Statistic for the Regression Model}

Because the regression model of Chapter~2 is a special case of the GMM model
of this chapter, it may be useful to know how $LR$ would look for that special case.
Because $\mathbf{x}_i = \mathbf{z}_i$ in the regression model, the (unrestricted) efficient GMM estimator
is OLS and $J(\hat{\boldsymbol{\delta}}(\hat{\mathbf{S}}^{-1}), \hat{\mathbf{S}}^{-1}) = 0$.
Therefore,
\begin{equation}
  LR = J(\tilde{\boldsymbol{\delta}}(\hat{\mathbf{S}}^{-1}),\, \hat{\mathbf{S}}^{-1}),
  \label{eq:3.7.3}
\end{equation}
where $\tilde{\boldsymbol{\delta}}(\hat{\mathbf{S}}^{-1})$ is the restricted efficient GMM estimator.
By Proposition~3.8, this statistic is asymptotically chi-squared and is numerically equal
to the Wald statistic if the null is linear.  As will be shown below, under conditional
homoskedasticity, this statistic can be written as the difference in the sum of squared
residuals normalized to the error variance.

%% ─── Variable Addition Test (optional) ──────────────────────────────────────
\subsection*{Variable Addition Test (optional)}

In the previous section, we considered a specification test based on the $C$ statistic
for the endogeneity of a subset, $\mathbf{x}_{i2}$, of instruments $\mathbf{x}_i$ while assuming that the
other instruments $\mathbf{x}_{i1}$ are predetermined.  Occasionally, we encounter a special case
where $\mathbf{z}_i = \mathbf{x}_{i1}$:
\begin{equation}
  y_i = \mathbf{x}_{i1}' \boldsymbol{\delta} + \varepsilon_i.
  \label{eq:3.7.4}
\end{equation}
A popular method to test whether the suspect instruments $\mathbf{x}_{i2}$ are predetermined is
to estimate the augmented equation
\begin{equation}
  y_i = \mathbf{x}_{i1}' \boldsymbol{\delta}
  + \mathbf{x}_{i2}' \boldsymbol{\alpha} + \varepsilon_i
  = \mathbf{x}_i' \boldsymbol{\gamma} + \varepsilon_i
  \quad \text{with } \boldsymbol{\gamma}
  = \begin{bmatrix} \boldsymbol{\delta} \\ \boldsymbol{\alpha} \end{bmatrix}
  \label{eq:3.7.5}
\end{equation}
and test the null $\mathrm{H}_0\colon \boldsymbol{\alpha} = \mathbf{0}$.  Testing is either by the Wald
statistic $W$ or by the $LR$ statistic, which is numerically equal to $W$.  This test is
sometimes called the \textbf{variable addition test}.  How is the test related to the $C$ test
of Proposition~3.7?

To calculate $LR$, we have to calculate two efficient GMM estimators of
$\boldsymbol{\gamma}$ with the same instrument set $\mathbf{x}_i$: one with the constraint
$\boldsymbol{\alpha} = \mathbf{0}$ and one without.  The unrestricted efficient GMM estimator is the OLS
estimator on the unconstrained equation \eqref{eq:3.7.5}.  The associated $J$ statistic is zero.
Let
\begin{equation}
  \hat{\mathbf{S}} \equiv \frac{1}{n} \sum_{i=1}^{n}
  e_i^2\, \mathbf{x}_i \mathbf{x}_i',
  \label{eq:3.7.6}
\end{equation}
where $e_i$ is the OLS residual from the unrestricted regression \eqref{eq:3.7.5}.  If we use
this estimate of $\mathbf{S}$, the restricted efficient GMM estimator of $\boldsymbol{\gamma}$ minimizes the $J$
function
\[
  J(\tilde{\boldsymbol{\gamma}},\, \hat{\mathbf{S}}^{-1})
  = n \cdot (\mathbf{s}_{xy} - \mathbf{S}_{xx}\tilde{\boldsymbol{\gamma}})'
  \hat{\mathbf{S}}^{-1}
  (\mathbf{s}_{xy} - \mathbf{S}_{xx}\tilde{\boldsymbol{\gamma}})
\]
subject to $\boldsymbol{\alpha} = \mathbf{0}$.  Clearly, the restricted estimator can be written as
\[
  \tilde{\boldsymbol{\gamma}}
  = \begin{bmatrix} \hat{\boldsymbol{\delta}}(\hat{\mathbf{S}}^{-1}) \\ \mathbf{0} \end{bmatrix},
\]
where $\hat{\boldsymbol{\delta}}(\hat{\mathbf{S}}^{-1})$ is the efficient GMM estimator with $\mathbf{x}_i$ as instruments on the restricted
equation~\eqref{eq:3.7.4}.  So
\begin{align*}
  LR &= n \cdot (\mathbf{s}_{xy} - \mathbf{S}_{xx}\tilde{\boldsymbol{\gamma}})'
    \hat{\mathbf{S}}^{-1}(\mathbf{s}_{xy} - \mathbf{S}_{xx}\tilde{\boldsymbol{\gamma}})
    \quad (\text{since } J = 0 \text{ for unrestricted GMM}) \\
  &= n \cdot (\mathbf{s}_{xy} - \mathbf{S}_{xx_1}\hat{\boldsymbol{\delta}}(\hat{\mathbf{S}}^{-1}))'
    \hat{\mathbf{S}}^{-1}
    (\mathbf{s}_{xy} - \mathbf{S}_{xx_1}\hat{\boldsymbol{\delta}}(\hat{\mathbf{S}}^{-1}))
    \quad (\text{since } \mathbf{S}_{xx}\tilde{\boldsymbol{\gamma}}
    = \mathbf{S}_{xx_1}\hat{\boldsymbol{\delta}}(\hat{\mathbf{S}}^{-1})).
\end{align*}
This is none other than Hansen's $J$ statistic for the restricted equation~\eqref{eq:3.7.4}
when $\mathbf{x}_i$ is the instrument vector.  This statistic, in turn, equals the $C$ statistic
because the $J_1$ in Proposition~3.7 is zero in the present case.  Therefore, all three
statistics, $W$ from the unrestricted regression, $LR$, and $C$, are numerically equal to
Hansen's $J$, provided that the same $\hat{\mathbf{S}}$ is used throughout.  That is, the variable
addition test is numerically equivalent to Hansen's test for overidentifying restrictions.

\medskip\noindent
\textsc{Questions for Review}

\begin{enumerate}
\item ($LR$ for the regression model)\quad Verify that, for the regression model where
  $\mathbf{z}_i = \mathbf{x}_i$,
  \[
    LR = \mathbf{y}'\mathbf{X}(n\hat{\mathbf{S}})^{-1}\mathbf{X}'\mathbf{y}
    - 2\mathbf{y}'\mathbf{X}(n\hat{\mathbf{S}})^{-1}(\mathbf{X}'\mathbf{X})
    \hat{\boldsymbol{\delta}}
    + \hat{\boldsymbol{\delta}}'(\mathbf{X}'\mathbf{X})(n\hat{\mathbf{S}})^{-1}
    (\mathbf{X}'\mathbf{X})\hat{\boldsymbol{\delta}}.
  \]

\item (Choice of $\hat{\mathbf{S}}$ in variable addition test)\quad Suppose you form
  $\bar{\mathbf{S}}$ from the residual from the restricted regression~\eqref{eq:3.7.4} and use it to
  form $W$, $LR$, and $C$.  Are they numerically equal?  Are they asymptotically chi-squared?
  \textbf{Hint:} Is this $\bar{\mathbf{S}}$ consistent under the null?
\end{enumerate}
